% =====================================================================
%  二、问题分析
%  写法约定：每问只说明三件事，即考虑到了什么、要建立什么样的模型、
%  为什么最终采用这一做法；推导细节留给第五章。
% =====================================================================
\section{问题分析}

四个子问题层层递进。问题一是单次交会定位的几何基础；问题二则将该几何模型进一步用作
检测点位置的选择依据，形成单源定位策略；问题三在源数未知的条件下，把单源定位能力嵌入“测向、收敛、
清除”的在线决策；问题四进一步考虑未知发射朝向，使定位、区域排除与航路规划适用于全向源和定向源混合的情形。

\subsection{问题一的分析}
本问不宜假设误差分布，而应利用题目给出的误差界，采用最坏情形模型：
每个检测点的示向度对应一个由“测量方向 \(\pm 1^\circ\)”张成的角楔。
多个检测点的角楔求交，即得干扰源的可能定位区域。由于角楔是凸集，其交仍为凸集
（有界时即为凸多边形），因此区域直径可转化为凸多边形顶点间的最远点对问题。求解时，需要先判断
交会区域是否有界，再检验以该直径为直径的圆能否覆盖定位区域。整体求解思路可概括为：由误差界构造单点角楔，
通过多检测点角楔交会确定凸多边形定位区域，进而将区域直径计算归结为顶点间的最远点对问题；
同时辅以有界性判定与圆覆盖检验。

\subsection{问题二的分析}
问题的关键在于第二个检测点无法由单次观测唯一确定：只掌握 $S_1$ 处的示向度，源点的
距离信息完全缺失，故既不能先定位再布点，也不存在一个位置在每次试验中都最好。
本文改以“布点后的定位效果之期望”反向衡量布点位置，沿用问题一中“定位区域直径”这一
判据，在中心射线与理想中心方位近似下对源位加权，并保留测向误差角楔。
在固定源位先验和忽略靶区边界截断的局部模型下，几何估价由第二检测点的相对位置决定。
于是以相对位置为决策变量、以 $E[D](a,b)$ 为价值函数进行二维寻优，选取低值区域作为补测候选。
该估价未对第二次实际接收事件建立完整概率模型，仅用于比较几何位置的平均效果。

\subsection{问题三的分析}

源数未知使“已发现目标全部清除”并不等于任务完成，因而本问需要同时解决定位清除、未知源搜索和停止判定。
对已发现目标，沿用问题一的角楔交会，并结合接收范围与靶区边界构造保守定位区域；再借助问题二的期望直径场，
选择有利于缩小区域的补测位置。对尚未发现的频道，则利用无信号观测逐步排除可能藏有源的区域。

两类任务均需移动和检测，若分别规划，容易造成重复绕行。因此，策略先由局部选点扩展为系统覆盖，再将搜索测站与
清除目标编入共同航路，进一步根据沿途观测调整测站、访问顺序及逐频道检测时机。为判断何时可以退出，采用覆盖网格
并计入单元尺寸余量，将离散检测结果转化为连续区域的排除依据。最终以完整退出时间衡量总成本，以清除比例检验任务完成程度，
通过同场景配对比较与机制消融分析各代改进的作用。

\subsection{问题四的分析}

定向源使无信号观测兼有距离过远和背向发射两种原因，因而需将未知状态扩展为源位与朝向的组合。
一次负观测只排除与其矛盾的组合，只有每个可能位置的全部朝向均被排除，才能判断无残留。
据此建立连续位置与朝向覆盖模型，以方格余量及朝向区间并集构造无源证书，并从不同方向布置测站；
对靶区边缘朝外发射的源，还需设置域外测站。

对已发现目标，保留角楔交会的保守定位区域，采用视线两侧的短基线补测，兼顾交会角与接收条件，
仅在可证明时利用双侧负观测截断远端；仍未达到定位精度时，以光学清除圆盘覆盖剩余区域。
进一步在覆盖约束下压缩测站布局，利用正负观测改善单方位目标估计，将搜索站与清除目标编入共同航路，
并计入访问顺序影响的后续检测成本。最终以连续覆盖保证完成、保守约束支持清除、联合规划降低耗时，
通过同场景配对实验检验效率改进。
